(1) 在固定 T 下最大化 \(f(x)=x^3/(e^x-1)\)。求导并整理得
\[
f'(x)=\frac{x^2[3(e^x-1)-xe^x]}{(e^x-1)^2}.
\]
非零驻值满足 \(3(1-e^{-x})=x\)。x=0 只是形式上该方程的端点，谱在 x=0 为零，不是峰值。

(2) 此方程在正实轴有唯一非零解 \(x_*\simeq2.821439\)，可用单调性或数值迭代确定。因此
\[
\omega_{\rm peak}=x_*k_BT/\hbar,
\]
频率峰线性随 T 增加，是 Wien 位移定律的角频率形式。

(3) 对全部角频率积分：
\[
u(T)=\int_0^\infty u_\omega\,d\omega
=\frac{(k_BT)^4}{\pi^2c^3\hbar^3}
\int_0^\infty\frac{x^3}{e^x-1}\,dx.
\]
利用积分 \(=\Gamma(4)\zeta(4)=\pi^4/15\)，得
\[
u(T)=\frac{\pi^2k_B^4}{15c^3\hbar^3}T^4.
\]
总能量密度依赖温度四次方。

(4) 将角频率改变量为真空波长 \(\lambda=2\pi c/\omega\)，谱满足 \(u_\lambda=u_\omega|d\omega/d\lambda|\)。Jacobian \(2\pi c/\lambda^2\) 会改变极大点，故波长谱峰满足 \(5(1-e^{-y})=y\)，其正根 \(y\simeq4.965114\)。因此频率峰的直接换算波长与波长谱峰不同，并无物理矛盾。
